\newcommand{\paramDecomp}{$\ell$}

\section{TFHE (RGSW)}
\label{sec:rgsw}
\ifnotes
\begin{itemize}
    \item Problem with relinearization (noise, keys, time etc.)
    \item Introduce (R)GSW as solution
    \item External product
    \item Trade-offs (size $\to$ communication overhead)
    \item Decomposition parameter \paramDecomp bounds
    \item HomExpand
\end{itemize}
\fi

The GSW ciphertext defined in \cite{cryptoeprint:2013/340}, now in the ring $R_q$. It works by changing bases. Given a base $B$ and decomposition parameter \paramDecomp, such that the message $\mu$ can be expressed with $\ell$ digits in base $B$, define the \emph{gadget vector} as $g = (1, B, \dots, B^{\ell - 1})^\intercal\bmod{t}$ and define the \emph{gadget matrix} $G$ as \eqref{eq:gadget-matrix}.

\begin{equation}
    \label{eq:gadget-matrix}
    G = I_2 \otimes g = 
    \begin{bmatrix}
        g & 0 \\
        0 & g
    \end{bmatrix}
\end{equation}

Then an \gls{rgsw} ciphertext $C$ encrypting a message $\mu \in [0, N)$ is defined by \eqref{eq:rgsw}, where each row of $Z$ is an \gls{rlwe} encryption of 0 under the secret key $s$ (effectively rows of public keys), $\mu$ is the message and $G$ is the gadget matrix (duh).

\begin{equation}
    \label{eq:rgsw}
    C = Z + \mu \cdot G
\end{equation}

% \textcolor{magenta}{I show in \autoref{sec:impl:rgsw} that each row $C_i$ can be constructed directly...}
While \eqref{eq:rgsw} is sufficient to construct $\mathsf{RGSW}_s(\mu)$, we can speed up the process by directly encrypting each row. (Note \gls{rgsw} can be viewed as a flat vector of $2\ell$ \gls{rlwe} ciphertexts where the first $\ell$ modify $a$ and the last $\ell$ modify $b$).

\noindent\textcolor{red}{\textbf{Maybe use $\mathsf{RLWE} = (b, a)$ like OpenFHE uses in e.g. $\to$\texttt{GetElements()}? In which case the top and bottom rows swap place.}}

\noindent\textcolor{magenta}{\textbf{Implementation details below, but ultimately not used}}

\newpage
\begin{corollary}
\label{corollary:rgsw:bottom}
For an \gls{rgsw} ciphertext $C$ encrypting a message $\mu$ under $s$, the $(\ell + i)$-th row $(0 \leq i < \ell)$ is equivalent to: 
\begin{equation}
    C_{\ell + i} = \mathsf{RLWE}_s(\mu B^i/\Delta)
\end{equation}

\begin{proof} 
    Follows from the definitions \eqref{eq:rgsw} and \eqref{eq:gadget-matrix}.
    \begin{align*}
        Z_i + \mu\cdot G_i &= \mathsf{RLWE}_s(0) + \mu \cdot(0, g_i) \\
            &= (a, a\cdot s + \epsilon) + (0, \mu\cdot B^i) \\
            &= (a, a\cdot s +\mu B^i + \epsilon) \\
            &= \mathsf{RLWE}_s(\mu B^i/\Delta)
    \end{align*}
\end{proof}
\end{corollary}

\noindent The same approach shows the first $\ell$ rows can be constructed as an \gls{rlwe} encryption $c_i = \mathsf{RLWE}_s(-s\mu B^i/\Delta)$ for $0 \leq i <\ell$ (\autoref{lemma:rgsw:top}). However, calculating $s\cdot \mu$ is not trivial in practice, and it's faster to simply modifying the $a$ component of a fresh $\mathsf{RLWE}(0)$ encryption. % TODO: Create from scratch
\autoref{alg:rgsw} creates $\mathsf{RGSW}_s(\mu)$ in $\mathcal{O}(\ell)$ steps (but can only be performed on the client-side as it requires $s$). 

%% Generate RGSW on client-side
%%%
%%% RGSW(m) on client-side
%%%
\begin{algorithm}[ht!]
\caption{Client-side RGSW generation}
\label{alg:rgsw}
\begin{algorithmic}[1]
    \State \textbf{Input:} Polynomial $\mu = \sum_{j} b_j X^j,\quad b_j\in[0,N)$
    \State \textbf{Input:} Decomposition parameter $\ell$
    \State \textbf{Output:} $\mathsf{RGSW}_s(\mu)$
    \State Construct a $2\ell\times 1$ vector $C$
    % \State let $z = 0$ % cache plaintext
    \For{$i = 0,\; i < \ell$}
        \State Let $\hat{\mu} = \mu B^i/\Delta = \frac{B^i}{\Delta}(\sum_{j} b_j X^j) \quad \textrm{note: coeff.}\;B^i\cdot b_j \bmod{q}$
        \State $\mathbf{C}_i = \mathsf{RLWE}_s(\hat{\mu})$
        \State $\mathbf{C}_{i+\ell} = \mathsf{RLWE}_s(0) + (0, \hat{\mu})$
    \EndFor
    \State\Return $\{\mathbf{C}_k\}^{2\ell - 1}_{k=0}$
\end{algorithmic}
\end{algorithm}


\noindent\textcolor{red}{Note: we focus on BGV, where $\Delta = 1$, which simplifies the algorithm even more.}

\ifnotes
\subsubsection{In BGV-RNS}
In \gls{rns}, each coefficient $b_i \in R_q$ is represented as a set of $k$ residues mod $q_i$ where $q = q_0q_1\cdots q_{k-1}$, such that $b_i$ is reconstructed using \gls{crt}.
\fi

\subsection{Parameter Choice}
The choice of basis $B$ and decomposition parameter $\ell$ is one and the same. For a given basis, the message polynomial $\mu$ is rewritten into $\log_B(\mu)$ base-$B$ digits. Therefore $\ell \geq \log_B(\mu)$, otherwise the RGSW does not encode all the digits, and $\ell \leq \log_B(\mu)$ because more $B$-digits would always be 0. The value of $\ell$ is then squeezed and must equal the number of base-$B$ digits that encodes all $\mu$ values allowed by the scheme.

%%------------------%%
%% EXTERNAL PRODUCT %%
%%------------------%%
\newpage
\subsection{External Product}
\ifnotes
\textcolor{magenta}{The purpose of using RGSW \emph{is} to get ``access'' to the external product, which allows homomorphic multiplication without the expensive relinearization step.}
\fi
First defined in \cite{cryptoeprint:2018/421}? The external product \eqref{eq:externalprod} is defined between an \gls{rlwe} and a \gls{rgsw} ciphertext, and outputs an \gls{rlwe} ciphertext without needing to use homomorphic multiplication (and thus relinearization). The external product uses a \emph{decomposition matrix} $G^{-1}$ that has the \emph{gadget property} \eqref{eq:gadget-property}, which effectively ``undoes'' the effect of the gadget matrix $G$ (with some error $\leq B/2$, assuming $x$ has binary elements).

\begin{align}
    \label{eq:gadget-property}
    G^{-1}(x)\cdot G &= x + e \pmod{q} \\
    \label{eq:externalprod}
    x \boxdot Y &= G^{-1}(x)\cdot Y
\end{align}


\subsection{BGV-RNS Optimizations}
Some key observations allow us to much more efficiently port the \gls{rgsw} operations to BGV-RNS. 

\begin{enumerate}
    \item RNS ciphertexts are decomposed by \gls{crt} using $\{ q_i \} \;\forall\; i$ as a basis.
    \item If we choose $g = (q_0, q_1, \dots, q_\ell)$ as the gadget vector, then we get gadget de/composition for ``free'' as it is already baked into the RNS structure
\end{enumerate}

\ifnotes
\textbf{No homomoprhic multiplication:}
\begin{align}
    x\boxdot Y &= G^{-1}(x)\cdot Y \\
    &= G^{-1}(x)\cdot(Z + \mu_y\cdot G) \\
    &= G^{-1}(x)\cdot Z + \mu_y \cdot G^{-1}(x)\cdot G \\
    &= \underbrace{G^{-1}(x)\cdot Z}_{A} + \underbrace{\mu_y \cdot (x + \epsilon)}_{B}
\end{align}

\textbf{Term A:}
???

\textbf{Term B:}
Since $x = (b, a) = (as + \Delta \mu_x + \epsilon')$ we get:
\begin{align}
   \mu_y\cdot &(as + \Delta \mu_x + \epsilon, a + \epsilon) \\
   &(as\cdot\mu_y + \Delta \mu_x\mu_y + \epsilon, a\mu_y + \epsilon) \\
   &((a\mu_y)\cdot s + \Delta (\mu_x\mu_y) + \epsilon, (a\mu_y) + \epsilon) \\
   &\mathsf{RLWE}_s(\mu_x\mu_y) + \epsilon
\end{align}

% As mentioned eariler, this does not use homomoprhic multiplication. \textcolor{red}{Show that this does not use (homomorphic) multiplication:} 

% \begin{align}
%     x \mod{B}
% \end{align}

\newpage
\fi
\ifnotes
\subsubsection{Gadget Decomposition}
The gadget encodes the message as digits in base $B$. To reconstruct the message $m$, we therefore need to find a decryption function such that:

\begin{equation}
      \mathsf{Dec}(\cdot): R_q \to R^{\ell} \quad s.t.\quad \langle \mathsf{Dec}(x), g\rangle = x \pmod{Q}
\end{equation}

Then $x\boxdot Y = \langle \mathsf{Dec}(x), Y\rangle$

% Encodes the message as digits in base $B$, which limits the noise to $B/2$ per level $\approx \mathcal{O}(B\log{B}) = \mathcal{O}(B\ell)$ instead of ...

% The RGSW really encodes $\mu$ as digits in $B$. In general, for each component of a vector, the gadget composition takes $B^i$. Let $v_B$ be $v$ in base $B$.

\begin{equation}
    (v_0, v_1, \dots, v_{\ell-1}) \cdot g = v_0 + v_1B + \dots + v_{\ell - 1} B^{\ell - 1} = v_B
\end{equation}

A general gadget matrix for a $k$-dimensional vector then becomes: 
\begin{equation}
    G = I_k \otimes g
\end{equation}

And the decomp:
\begin{align}
    G^{-1}(x)\cdot G &= G^{-1}(x)\cdot(I_k\otimes g) = x + \epsilon \\
    G^{-1}(x) &= (I_k \otimes g)^{-1}(x + \epsilon)
\end{align}

In our case, each component of the vector $v$ is an \gls{rlwe} encryption $(b,a)$, so $k = 2$. 
\begin{equation}
    G^{-1}(x) = 
    \begin{pmatrix}
        g & 0 \\
        0 & g
    \end{pmatrix}
    ^{-1}
    (a + \epsilon, b + \epsilon) = 
    \frac{1}{}
\end{equation}


From \label{eq:rgsw} we can see that we are gadgeting the message $\mu$. In other words, assuming we are encrypting bits, $\mu$ is encoded into base $B$ digits in each component of our ciphertext(s). 
\begin{equation}
    Z + \mu B
\end{equation}

In one dimension, the gadget property \eqref{eq:gadget-property} for a decomposition vector $g^{-1}$ undoes the gadget:
\begin{equation}
    (v_0 + v_1B + \dots + v_{\ell - 1} B^{\ell - 1}) \cdot g^{-1} = (v_0, v_1, \dots , v_{\ell - 1})
\end{equation}

\subsubsection{Constructing Decomp Matrix}
????????

\subsubsection{Placing}
Let $u \in \{0,1\}$.

\begin{equation}
    \mathsf{RLWE}_s(v) \boxdot \mathsf{RGSW}_s(u) = 
\end{equation}

Alternative: just show 1) it exists and 2) this matrix has the property
\fi

\newpage

\subsection{Internal Product}
RGSW $\times$ RGSW

\begin{equation}
    A \boxtimes B = 
    \begin{bmatrix}
        A \boxdot b_0 \\
        A \boxdot b_1 \\
        \vdots \\
        A \boxdot b_\ell
    \end{bmatrix}
\end{equation}

\newpage

%%-----------------------%%
%% Homomorphic Expansion %%
%%-----------------------%%
\subsection{RLWE-to-RGSW Expansion from ORAM}
\ifnotes
\begin{itemize}
    \item We want to use RGSW to avoid relinearization
    \item We want to send packed ciphertexts to save on communication
    \item Therefore, we send RLWE and convert to RGSW on the server 
    \item (Requires a single RGSW encryption $\mathsf{RGSW}(-s)$ to be sent as well)
    \item 1 RLWE for the index bits = size $\ell$
    \item $\ell$ RGSW's for the index bits = size $2\ell \times 2$ or smth
    \item Expanding = size $2\ell$ or something
\end{itemize}
\fi
As established in \autoref{sec:rgsw}, \gls{rgsw} ciphertexts are enticing because the external product allows us to perform homomorphic multiplication without relinearization. However, the size of \gls{rgsw} makes the communication costs very high in a communication protocol. We wish to strike a balance between the cheap communication cost of packed \gls{rlwe} and the fast multiplication of \gls{rgsw}. The solution is to send \gls{rlwe} ciphertexts from the client, and \emph{expand} them homomorphically on the server.

Solution from \autocite{chillotti:onionring:2019}. \textbf{HomExpand:} Expands an \gls{rlwe} ciphertext into an \gls{rgsw} ciphertext.

This expansion is described in Algorithm 4 of \cite{chillotti:onionring:2019}, which builds on Algorithms 3 and 8 from the same paper. The algorithms are listed below (in slightly different notation) as \autoref{alg:ORAM:4}, \autoref{alg:ORAM:3}, and \autoref{alg:ORAM:8}, respectively.

%%%
%%% ORAM ALGORITHM 4 - HomExpand (RLWEtoRGSW)
%%%
\algnewcommand\algorithmicforeach{\textbf{for each}}
\algdef{S}[FOR]{ForEach}[1]{\algorithmicforeach\ #1\ \algorithmicdo}

%%%
%%% ORAM Algorithm 4 - HomExpand (RLWEtoRGSW)
%%%
\begin{algorithm}[ht!]
\caption{HomExpand}
\label{alg:ORAM:4}
\begin{algorithmic}[1]
    \State \textbf{Input:} $\mathbf{c}_k$ = RLWE$(\sum_{i=0}^{n-1}\frac{b_i}{nB^{k+1}}),\; b_i\in\{0,1\},\; 0 \leq k < \ell$
    %, \Statex $\quad\quad\quad\;\; k=0,1,\dots,\ell - 1$
    \State \textbf{Input:} $A = \mathrm{RGSW}(-s)$
    \State \textbf{Output:} $\mathbf{C}_i = \mathrm{RGSW}(b_i),\; 0 \leq i < n$
    \ForEach{$k$} % in $0, \dots, \ell - 1$}
        \State $\mathbf{c'}_k:= \mathsf{expandRlwe}(c_k)$
    \EndFor
    \For{$i = 0,\; i < n$}
        \State Initialize $\mathbf{C}_i$ as an empty $2\ell \times 2$ matrix of polynomials
        \For{$k = 0; k < \ell$}
            \State $\mathbf{C}_i[k] = A \boxdot \mathbf{c'}_k[i]$
            \State $\mathbf{C}_i[k + \ell] = \mathbf{c'}_k[i]$
        \EndFor
    \EndFor
    \State\Return $\{\mathbf{C}_j\}^{n-1}_{j=0}$
\end{algorithmic}
\end{algorithm}

\autoref{alg:ORAM:4} takes as input $\ell$ RLWE ciphertexts, which in \gls{spar} correspond to the $\ell$ encrypted bits of the index \textcolor{red}{$a_j$}. It also takes an RGSW ciphertext of the encrypted secret key \textcolor{red}{(which key, automorphism? Sent by the client or generated by server?)}.

It begins by running a subroutine $\mathsf{expandRlwe}$ (\autoref{alg:ORAM:3}) which converts each of the $\ell$ RLWE ciphertexts into a matrix, like a \emph{partial} RGSW. Then these partial RGSW's are combined into a single (complete) RGSW.

%%%
%%% ORAM ALGORITHM 3 - RLWE Expansion subroutine ExpandRLWE
%%%
\input{Algorithms/oram-3}

This subroutine converts an RLWE ciphertext with $\ell$ slots into $\ell$ ciphertexts, each encoding a unique slot of the input. \textcolor{red}{Subs} is a function that... (EvalAutomorphism!)

Their method is fast for TFHE, but for our use --- case where $n << N$ --- it's faster to use the method described in \cite{cryptoeprint:2018/244} to repeatedly ``rotate'' the ciphertext $\ell$ times as in \autoref{alg:expand:bgv}. The FastRotate operation is supported by the OpenFHE as $\mathsf{EvalFastRotate}$.

Note that the components being scaled by $n$ is a side-effect of their implementation, and not required when using FastRotate, so the output is a list of $c_i =\mathsf{RLWE}(b_i)$.

%%
%% ORAM ALGORITHM 3 - MODIFIED
%%
%%%
%%% ORAM Algorithm 3 - RLWE Expansion subroutine ExpandRLWE
%%%
\begin{algorithm}[ht!]
\caption{expandRlwe}
\label{alg:expand:bgv}
\begin{algorithmic}[1]
    \State \textbf{Input:} $\mathbf{c} = \mathrm{RLWE}(\sum_{i=0}^{n-1} b_i X^i),\;b_i\in\{0,1\}$
    \State \textbf{Output:} $\mathbf{c}_i = \mathrm{RLWE}(b_i),\; 0 \leq i < n$
    \State $\mathbf{c}_0 := \mathbf{c}$
    \For{$i = 0,\; i \leq \log{n}$}
        \State do something...
        % \State $k = n/2^{i - 1} + 1$
        % \For{$b = 0; b < 2^i$}
        %     \State $\mathbf{c}_{2b} = c_b + \mathrm{Subs}(c_b, k)$
        %     \State $\mathbf{c}_{2b + 1} = c_b - \mathrm{Subs}(c_b, k)$
        % \EndFor
    \EndFor
    \State\Return $\{\mathbf{c}_j\}^{n-1}_{j=0}$
\end{algorithmic}
\end{algorithm}



\ifnotes
\subsubsection{Substitution}
Finally, the substitution is a bit more complex... Uses an automorphism key to rotate? Or something?

%%% -----------------------
%%% ORAM ALGORITHM 8 - Subs
%%% -----------------------
%%%
%%% ORAM Algorithm 8 - Subs
%%%
\begin{algorithm}
\caption{RLWE Substitution Algorithm $\mathsf{Subs}(c, k)$}
\label{alg:ORAM:8}
\begin{algorithmic}[1] % [1] adds line numbering every line
    \State \textbf{Input:} $\mathbf{c} = \mathrm{RLWE}_s(\mu(X))$; odd integer $k$
    \State \textbf{Input:} $KS_{s(X^k)\to s}[j]$ for $j = 1 : \ell_{ks}$
    \State \textbf{Output:} $\mathbf{c'} = \mathrm{RLWE}_s(\mu(X^k))$
    \State $\mathbf{c} = (c_0, c_1)$
    \State Let $\mathbf{c'} = (c'_0, c'_1) = (c_0(X^k), c_1(X^k))$
    \State Decompose $c'_0$ as $\sum_{j=1}^{\ell_{ks}} c'_0[j]/B^j_{ks} + \epsilon_{dec}$ with
    \Statex $\quad||c'_0[j]||_\infty \leq B_{ks}/2$ and $||\epsilon_{dec}||_\infty\leq 1/(2B^\ell_{ks})$
    \State\Return $\mathbf{c'} = (0, c_1) - \sum_{j=1}^{\ell_{ks}} c_0[j] \cdot KS_{s(X^k)\to s}[j]$
\end{algorithmic}
\end{algorithm}

\newpage % TODO: Fix floating issues...
\clearpage
\newpage
\fi


% \begin{algorithm}
% \caption{HomExpand}
% \label{alg:HomExpand}
% \begin{algorithmic}[1]
%     \State \textbf{Input:}  A set
%            $\mathcal{S} = \bigl\{c_j=(a_j,b_j)\bigr\}_{j=1}^{2\ell}$
%            of \gls{rlwe} ciphertexts under an auxiliary key $s'$, where $c_j$
%            encrypts $\bigl\lfloor q/\beta^{(j-1)\bmod\ell}\bigr\rfloor\cdot m$
%     \State \textbf{Output:} An \gls{rgsw} ciphertext
%            $\mathbf{C}\in\mathcal{R}_q^{2\ell\times 2}$
%     \State Initialise $\mathbf{C}\leftarrow\mathbf{0}^{2\ell\times 2}$,
%            $\;j\leftarrow 1$
%     \For{$c=(a,b)\in\mathcal{S}$}
%         \State $(a,b)\leftarrow\mathsf{KS}_{s'\to s}(a,b)$
%                \Comment{Re-encrypt under target key $s$}
%         \State $(a,b)\leftarrow\mathsf{ModSwitch}_{\,q'\to q}(a,b)$
%                \Comment{Reduce noise via modulus switching}
%         \State $\mathbf{C}[j,\,:\,]\leftarrow(a,\,b)$
%                \Comment{Assign as $j$-th row of RGSW matrix}
%         \State $j\leftarrow j+1$
%     \EndFor
%     \State \Return $\mathbf{C}$
% \end{algorithmic}
% \end{algorithm}


%%----------------%%
%%                %%
%%----------------%%
\subsection{Homomorphic Field Trace}
New paper!


%%----------------%%
%% IMPLEMENTATION %%
%%----------------%%
\newpage % Covered by lemma 2.1
\section{Implementation Details}

\subsection{RGSW}
\label{sec:impl:rgsw}
In practice, it would be convenient to construct each row directly, rather than first constructing both matrices in $\eqref{eq:rgsw}$ and combining them.

Recall that an RLWE ciphertext is a vector $\mathsf{RLWE}(\mu) = (a, b)$ where $b=a\cdot s + \mu + \epsilon$ \textit{(since $\Delta = 1$ in BGV)}. We can view \gls{rgsw} as a flat vector of $2\ell$ \gls{rlwe} ciphertexts, or a $2\ell \times 2$ matrix with an $a$ column and $b$ column. In the latter representation, the first $\ell$ rows of $C$ are expressed as:
\begin{equation*}
    C_i =
    (a,\; b_{\mu=0}) +
    \mu \cdot (1/B^{i},\; 0)
    =
    (a + \mu/B^{i},\; a\cdot s + \epsilon)
\end{equation*}

As each row of an \gls{rgsw} ciphertext is a valid \gls{rlwe} ciphertext, we can find a new message $\hat{\mu}_i$ such that $\mathsf{RLWE}(\hat{\mu_i}) = (\hat{a}, \hat{a}\cdot s + \Delta \cdot\hat{\mu}_i + \epsilon) = C_i$. Starting with the top half, let $\hat{a} = a + \mu\cdot B^{i}$ so that the first component matches.

\begin{align}
    b &= (a + \mu\cdot B^{i})\cdot s + \Delta \cdot \hat{\mu}_i + \epsilon \\
      &= a\cdot s + s\cdot \mu \cdot B^{i} + \Delta \cdot \hat{\mu}_i + \epsilon \\
      &= a\cdot s + (s\mu B^{i} + \Delta{\hat{\mu}}_i) + \epsilon
\end{align}

The middle term must equal 0 to match $C_i$, which leads us to an expression \eqref{eq:rgsw:top} for $\hat{\mu}_i$.
\begin{align} 
    0 &= s\cdot \mu \cdot B^{i} + \Delta\cdot \hat{\mu}_i \\
    \hat{\mu}_i &= -s\cdot\mu\cdot B^{i} \cdot 1/\Delta \\
    \label{eq:rgsw:top}
    \hat{\mu}_i &= -s\cdot\mu\cdot B^{i}
\end{align}

$\therefore$ we can construct row $i$ (0-indexed) of $\mathsf{RGSW}_s(\mu)$ as $\mathsf{RLWE}_s(-s\cdot\mu\cdot B^{i})$ for $0 \leq i < \ell$.Further, the RLWE plaintext $\mu$ in mod $t$ so we can also use $\mathsf{RLWE}_s(t - s\mu B^{i})$.

\begin{align*}
    \mathsf{RLWE}_s(-s\mu B^{i}) &= (a, a\cdot s - s\mu B^{i} + \epsilon) \pmod{t} \\
    \mathsf{RLWE}_s(t - s\mu B^{i}) &= (a, a\cdot s + (t - s\mu B^{i}) + \epsilon) \pmod{t} \\
    \mathsf{RLWE}_s(-s\mu B^{i})  &\equiv \mathsf{RLWE}_s(t - s\mu B^{i})   \pmod{t}
\end{align*}

For the bottom $\ell$ rows, we can skip Z entirely since, as \autoref{lemma:rgsw:bottom} says: 
$$\mathsf{RLWE}_s(0) + (0, \mu B^{i}) = (a,\; a\cdot s + \mu B^{i} + \epsilon) = \mathsf{RLWE}_s(\mu B^{i})$$

% \newpage
% Maybe: present as lemmas for top/bottom rows?
% \begin{lemma}
% \label{lemma:rgsw:bottom}
% For an \gls{rgsw} ciphertext $C$ encrypting a message $\mu$ under $s$, the $(\ell + i)$-th row ($0 \leq i < \ell$) is given by: 
% \begin{equation}
%     C_{\ell + i} = \mathsf{RLWE}_s(\mu B^i/\Delta)
% \end{equation}

% \begin{proof} 
%     Follows from the definitions \eqref{eq:rgsw} and \eqref{eq:gadget-matrix}.
%     \begin{align*}
%         Z_i + \mu\cdot G_i &= \mathsf{RLWE}_s(0) + \mu \cdot(0, g_i) \\
%             &= (a, a\cdot s + \epsilon) + (0, \mu\cdot B^i) \\
%             &= (a, a\cdot s +\mu B^i + \epsilon) \\
%             &= \mathsf{RLWE}_s(\mu B^i/\Delta)
%     \end{align*}
% \end{proof}
% \end{lemma}

\subsubsection{Appendix?}
\begin{lemma}
\label{lemma:rgsw:top}
The $i$-th row $C_i$ of $\mathsf{RGSW}_s(\mu)$ is:
\begin{equation}
    C_i = \mathsf{RLWE}_s(-s\mu B^{i}/\Delta)
\end{equation}

\begin{proof}
Let $\hat{a} = a + \mu B^i$, and consider $\mathsf{RLWE}_s(\hat{\mu})$.
\begin{align*}
    \mathsf{RLWE}(\hat{\mu}) &= (\hat{a}, \hat{a}\cdot s + \Delta \hat{\mu} + \epsilon) \\
        &=(a + \mu B^i, (a + \mu B^i)\cdot s + \Delta \hat{\mu} + \epsilon) \\
        &=(a + \mu B^i, a\cdot s + \mu B^i\cdot s + \Delta\hat{\mu} + \epsilon)
\end{align*}

Assume $C_i = \mathsf{RLWE}(\hat{\mu})$, then:
\begin{align*}
    s\mu B^i + \Delta\hat{\mu} &= 0 \\
    \therefore \hat{\mu} &= -s\mu B^i /\Delta
\end{align*}
\end{proof}
\end{lemma}

% \subsubsection{Prove thing}
% $enc(-s * msg * B^i) = Enc(0) - Enc(msg * B^i)$

% \begin{align}
%     \mathsf{RLWE}_s(0) - \rlwe{\mu B^i} &= (a, as + (0 - \mu B^i) + \epsilon) \\
%         &= (a, as - \mu B^i + \epsilon)
% \end{align}


\subsubsection{RNS optimizations}
OpenFHE is built on a \gls{rns} where each polynomial is represented as residues of smaller (typically 48 bits for AVX512 or 60 bits by default) primes. The polynomials are equivalent to their non-RNS counter-parts by the Chinese Remainder Theorem.

The base should be 60 (or 48) aka. \texttt{NativeSize} (?) in OpenFHE, then the gadget decomposition param becomes (lower bound)...

% \subsubsection{Implementation Details (Shorter)}
% In practice we wish to construct $C$ directly, rather than follow \eqref{eq:rgsw} directly (by first constructing $Z$ and $G$). An \gls{rgsw} can be viewed as a vector of $2\ell$ \gls{rlwe} ciphertexts. Note also $G$ is a sparse matrix represented uniquely by $g$. 

% \begin{equation}
%     C = (\mathbf{c_0}, \mathbf{c_1}) = \mathsf{RLWE}(0) + \mu\cdot g
% \end{equation}

% \begin{equation}
%     C_{i} = 
%     \begin{cases}
%         (a, b_0 - \mu \cdot B^i) = (a, b_{- \mu B^{i}}), & 0 \leq i < \ell \\
%         (a - \mu \cdot B^{i-\ell}, b), & \ell \leq i < 2\ell \\
%     \end{cases}
% \end{equation}


\newpage
\subsection{External Product for BGV-RNS}
\label{sec:impl:extprod}
Expand the condition \eqref{eq:externalprod:G}:

\newcommand{\rlwe}[1]{\mathsf{RLWE}(#1)}

% \begin{equation}
%     G^{-1}(x)\cdot
%     \begin{pmatrix}
%         a & a\cdot s -s\cdot\mu/B + \varepsilon \\
%         \vdots \\
%         a & a\cdot s -s\cdot\mu/B^\ell + \varepsilon \\
%         a & a\cdot s + \mu/B + \varepsilon \\
%         \vdots \\
%         a & a\cdot s + \mu/B^\ell + \varepsilon \\
%     \end{pmatrix}
%     =
%     \begin{pmatrix}
%         a, a\cdot s + \mu_a + 
%     \end{pmatrix}
% \end{equation}

\begin{equation}
    G^{-1}(x)\cdot 
    \begin{pmatrix}
        g & 0 \\
        0 & g
    \end{pmatrix} 
    = (a, \; a\cdot s + \mu)
    + \epsilon \pmod{q}
\end{equation}

% If we let $G^{-1}$ be a vector of size $2\ell$, then the first $\ell$ components build the public component $a$ ...
Due to the 0's, the top $\ell$ rows correspond to the public component $a$, while the bottom $\ell$ rows correspond to the encrypted message $b$. So $G^{-1}$ is a vector with $2\ell$ components, which we can split into two distinct vectors denoted $G^{-1}_\mathrm{top}$ and $G^{-1}_\mathrm{bot}$ for now.

\begin{align}
    G^{-1}_\mathrm{top}(x) \cdot g &= \sum_{i=1}^{l}u_i/B^i \approx a + \epsilon \pmod{q} \\
    G^{-1}_\mathrm{bot}(x) \cdot g &= \sum_{i=1}^{l}v_i/B^i \approx a\cdot s + \mu_x + \epsilon \pmod{q}
\end{align}

\begin{equation}
    G^{-1}_\mathrm{RNS}(a)_j = a \bmod \alpha_j \in R_{\alpha_j},
    \qquad
    g_j = \hat{\alpha}_j = \frac{q}{\alpha_j}
\end{equation}
\begin{equation}
    G^{-1}_\mathrm{RNS}(a) \cdot g 
    = \sum_{j=0}^{\ell-1} (a \bmod \alpha_j) \cdot \hat{\alpha}_j 
    \equiv a \pmod{q}
\end{equation}



% Top row can be written as:
% \begin{equation}
%     G^{-1}_\mathrm{top}(x)\cdot
%     \begin{pmatrix}
%         1/B \\ \vdots \\ 1/B^{\ell}
%     \end{pmatrix}
%     = \sum_{i=1}^{\ell}\frac{[G^{-1}_\mathrm{top}(x)]_i}{B^i} 
%     \approx a \pmod{x}
% \end{equation}

% Let $g^{-1} = B + BX + \dots + B^\ell X^\ell = (B, \dots, B^\ell)$. Clearly, if  $G^{-1}_{top}(x) = [a] \cdot g^{-1}$ works (as $g^{-1}_i \cdot g_i = 1$ for $0 <i \leq \ell$, where $g_i$ denotes the $i$-th coefficient in the polynomial $g$.) This is fine because $G^{-1}(x)$ takes $x = (a,b)$ as input, so $a$ is easily accessible. The same argument can be applied to see that $b + g^{-1}$ is a valid choice for $G^{-1}_{bot}(x)$ as well.

% \begin{equation}
%     (b + g^{-1}) \cdot g = \sum_{i = 1}^{\ell}\frac{b_i \cdot B^i}{B^i} = b \pmod{x}
% \end{equation}

% % Let $g^{-1} = (B, \dots, B^\ell) \in \mathbb{Z}_q^{1\times\ell}$ be the row vector of gadget inverses, satisfying $g^{-1}_i + g_i = 1$ for $0 < i \leq \ell$ (where $g_i$ denotes the $i$-th entry of $g$).
% % Then $G^{-1}_\mathrm{top}(x) = a \cdot g^{-1}$ is a valid choice, since $G^{-1}(x)$ takes $x = (a, b)$ as input, so $a$ is readily accessible. The same argument gives $G^{-1}_\mathrm{bot}(x) = b \cdot g^{-1}$, yielding:

So in the code, we consider $G^{-1}(x)$ as two separate vectors/polynomials $a\cdot g^{-1}$ and $b \cdot g^{-1}$ and then calculate the external product \eqref{eq:externalprod} as: 

\begin{align}
    x \boxdot Y &= (a, b) \\
    \mathrm{where}\;\,  a &= (x_1\cdot g^{-1}) \cdot Y \\
                        b &= (x_2\cdot g^{-1}) \cdot Y
\end{align}

\begin{equation}
    a = \sum_{i=0}^n c_i\cdot x^i \quad\implies\quad a\cdot b = \sum 
\end{equation}

% \begin{equation}
%     g^{-1}(x)\cdot 
%     \begin{pmatrix}
%         1/B      & 0        \\
%         \vdots   & \vdots   \\
%         1/B^\ell & 0        \\
%         0        & 1/B      \\
%         \vdots   & \vdots   \\
%         0        & 1/B^\ell
%     \end{pmatrix} 
%     = (a, \; a\cdot s + \mu)
%     + \epsilon \pmod{q}
% \end{equation}

% Examine one of the bottom $\ell$ rows, which clearly build the b component:

% \begin{equation}
%     g^{-1}_{\ell + 1}(x)\cdot 1/B + \dots + g^{-1}_{2\ell}(x) \cdot 1/B^\ell
%     = \sum_{i = \ell + 1}^{2\ell}{\frac{g^{-1}_{i}(x)}{B^{i}}}
% \end{equation}

% \begin{equation}
%     \sum_{i = \ell + 1}^{2\ell}{\frac{g^{-1}_{i}(x)}{B^{i}}} \approx a\cdot s + \mu
% \end{equation}


% Let $g^{-1} = x \mod{B}$

% \begin{equation}
%     g^{-1}(a)\cdot 
%     \begin{pmatrix}
%         1/B \\
%         \vdots \\ 
%         1/B^{\ell}
%     \end{pmatrix} 
%     = a + \epsilon \pmod{q}
% \end{equation}

% \begin{align*}
%     G^{-1}(a)\cdot G &= a + e &\mod{q} \\
%     g^{-1}(a)\cdot (1/B, \dots, 1/B^{\ell})^\intercal\bmod{t} &= a + e &\mod{q} \\
%     g^{-1}(a) &\approx a\cdot B^{i} &\mod{q} \\
%     (u, v) &\approx (a\cdot B^{i}, (a\cdot s + \mu)\cdot B^{i}) &\mod{q} \\
% \end{align*}

% \newpage
% In the BGV-RNS setting, the ciphertext modulus $q = \prod_{j=0}^{\ell-1} \alpha_j$
% is factored into $\ell$ digit moduli. The gadget vector is
% $g = (\hat{\alpha}_0, \ldots, \hat{\alpha}_{\ell-1})^\top$ where
% $\hat{\alpha}_j = q/\alpha_j$, and the decomposition is:
% \begin{equation}
%     [G^{-1}(a)]_j = a \bmod \alpha_j
% \end{equation}
% This is a projection onto the $j$-th RNS digit — no per-coefficient
% arithmetic — and the reconstruction $\sum_j (a \bmod \alpha_j)\cdot\hat{\alpha}_j = a$
% holds by CRT. The external product is then:
% \begin{equation}
%     \mathbf{x} \mathbin{\boxdot} Y
%     = G^{-1}_\mathrm{RNS}(\mathbf{x}) \cdot Y
% \end{equation}
% where each digit $[G^{-1}(\mathbf{x})]_j \in R_{\alpha_j}^{1\times 2}$
% is computed by truncating the RNS representation of $\mathbf{x}$ to
% the $j$-th limb group.

% \subsubsection{Implementation Details (AI)}

% The gadget $G = \mathrm{diag}(g, g) \in \mathbb{Z}_q^{2\ell \times 2}$ with
% $g = (1, B, \ldots, B^{\ell-1})^\top$ is chosen so that
% $G^{-1}(\mathbf{x}) \cdot G \approx \mathbf{x} \pmod{q}$.
% With $\mathbf{x} = (a, b) \in R_q^2$ and writing
% $G^{-1}(\mathbf{x}) = (u_0, \ldots, u_{\ell-1},\; v_0, \ldots, v_{\ell-1})
% \in R_q^{1 \times 2\ell}$, the block structure of $G$ decouples into two conditions:
% \begin{equation}
%     \sum_{i=0}^{\ell-1} u_i \cdot B^i = a,
%     \qquad
%     \sum_{i=0}^{\ell-1} v_i \cdot B^i = b
%     \pmod{q}
%     \label{eq:decomp-conditions}
% \end{equation}
% These are satisfied by the \emph{digit decomposition}: apply base-$B$ expansion
% to every coefficient of $a$ simultaneously.
% Define the decomposition $\mathrm{Decomp}_B : R_q \to R_q^{\ell}$ by
% \begin{equation}
%     r_0 := a, \qquad
%     u_i := r_i \bmod B, \qquad
%     r_{i+1} := \bigl\lfloor r_i / B \bigr\rfloor
%     \qquad (0 \le i < \ell)
% \end{equation}
% where $\bmod$ and $\lfloor\cdot/B\rfloor$ act \emph{coefficient-wise} on the
% polynomial. By construction:
% \begin{equation}
%     \sum_{i=0}^{\ell-1} u_i \cdot B^i = a \pmod{q}
%     \label{eq:decomp-exact}
% \end{equation}
% since each $u_i$ holds the $i$-th base-$B$ digit of every coefficient
% and each digit satisfies $u_i \in (-B/2,\, B/2]$ (small).
% Applying the same procedure to $b$ gives $\mathrm{Decomp}_B(b)$, and the
% external product \eqref{eq:externalprod} is then:
% \begin{equation}
%     \mathbf{x} \mathbin{\boxdot} Y
%     \;=\; G^{-1}(\mathbf{x}) \cdot Y
%     \;=\;
%     \bigl(\,\mathrm{Decomp}_B(a) \;\big|\; \mathrm{Decomp}_B(b)\,\bigr)
%     \cdot Y
%     \label{eq:extprod-impl}
% \end{equation}
% which in code is: decompose $a$ and $b$ into $\ell$ small-coefficient
% polynomials each, concatenate into a $1 \times 2\ell$ row, and multiply
% into the RGSW matrix $Y \in R_q^{2\ell \times 2}$.